\begin{tcolorbox}[colback=propbox, colframe=propborder, arc=2pt,
  left=6pt, right=6pt, top=4pt, bottom=4pt,
  title={\small\textbf{Definition (Logical Connective)}},
  fonttitle=\small\bfseries]
\label{def:logical-connective}
The standard primitive \emph{logical connectives} are:
\begin{center}
\renewcommand{\arraystretch}{1.25}
\begin{tabular}{c c c c l}
\toprule
\textbf{Symbol} & \textbf{Name} & \textbf{Arity} & \textbf{Formation pattern} & \textbf{Reading} \\
\midrule
$\neg$ & Negation & unary & $\neg\varphi$ & not $\varphi$ \\
$\land$ & Conjunction & binary & $(\varphi\land\psi)$ & $\varphi$ and $\psi$ \\
$\lor$ & Disjunction & binary & $(\varphi\lor\psi)$ & $\varphi$ or $\psi$ \\
$\to$ & Conditional & binary & $(\varphi\to\psi)$ & if $\varphi$, then $\psi$ \\
$\leftrightarrow$ & Biconditional & binary & $(\varphi\leftrightarrow\psi)$ & $\varphi$ iff $\psi$ \\
\bottomrule
\end{tabular}
\end{center}
\end{tcolorbox}

\begin{remark*}[Interpretation]
Arity records how many formulas a connective uses to form a new formula. The
truth behavior of these connectives belongs to the semantics section, not to
the syntactic definition.
\end{remark*}

\begin{dependencies}
\begin{itemize}
  \item \hyperref[def:propositional-language]{Propositional Language}
\end{itemize}
\end{dependencies}
